\section{Introduction}
\label{sec:intro}
Two of the largest low-volatility equity funds are built on opposite
philosophies. The iShares MSCI USA Min Vol Factor ETF, USMV, tracks an index
defined by a constrained quadratic programme over a commercial multi-factor
risk model \citep{ishares2011usmv497k}. The Invesco S\&P 500 Low Volatility
ETF, SPLV, tracks an index defined by two sentences: take the hundred least
volatile members of the parent, weight them by the inverse of their volatility
\citep{invesco2024splv497k}. Both providers publish their rules. Only one of
them publishes rules a reader can act on.

That asymmetry is the subject of this paper. A rank-and-weight rule is complete
in the sense that two people implementing it from the document will produce the
same portfolio. A constrained optimisation is not, because the constraints are
written in the vocabulary of the provider's own risk model. The weight cap,
the sector band and the turnover limit survive translation. The style-factor
band does not: it is stated relative to standardised exposures to factors the
provider names but does not define, and a replicator with a different factor
set has to decide what those names map onto. The decision is invisible in the
methodology document and consequential in the result.

This is a different problem from index tracking as that literature usually
poses it. There the index is taken as given and the task is to approximate its
return series with a manageable subset of its constituents, whether by
optimisation, by statistical factorisation or by learned models; recent
treatments include \citet{cesarone2025benchmark}, and \citet{dhingra2026index}
survey the field. We are not approximating a return series. We are
reconstructing the rule that generates it, from the document that states the
rule, and asking what that document leaves undetermined. Tracking error appears
here as the measure of how well the reconstruction succeeded, not as an
objective anything was optimised against.

We quantify how consequential. Using a weekly factor risk model built entirely
from free data, we rebuild both indices over the funds' own lifetimes and
measure each replica against the fund that tracks the index it targets. The
SPLV replica achieves 1.62\% annualised tracking error with a correlation of
0.9939, close enough that the remaining difference is mostly the
parent-universe proxy. The USMV replica achieves 5.83\%. Closing that gap is
where the paper's content lies, and the usual suspect turns out to be the wrong
one.

The usual suspect is the covariance matrix. If a commercial risk model is the
thing a free-data replication cannot reproduce, then a better covariance
estimate should narrow the gap, and the literature offers a well-stocked menu
of them. We run four: the sample matrix, a two-half-life structural factor
model in the Barra idiom, constant-correlation shrinkage after
\citet{ledoit2003improved}, and a principal-component model whose factor count
is chosen by the \citet{bai2002determining} criteria rather than assumed. Under
one specification they track the fund within 26 basis points of one another. That is the ordering \citet{dom2025beyond} report for constrained
portfolios in general, and \citet{jagannathan2003risk} supply the reason:
binding weight constraints act as a form of shrinkage, so a portfolio that is
already constrained has less to gain from a better estimate of $\Sig$. Our
contribution on this point is a confirmation in a specific published
constraint set, not a new result.

The specification choice is worth far more. The methodology exempts one
proprietary volatility factor from the style bands, which is the entire
mechanism by which a minimum-volatility index is permitted to be a
low-volatility index. A public-data model has two candidates for what that
factor names, beta and residual volatility, and the document that governs the
index in question does not say which. Exempting residual volatility alone gives
5.83\%. Exempting beta as well gives 3.73\%, with replica volatility of 13.97\%
against the fund's 13.45\% and a beta to the fund of 1.001. One sentence in a
constraint document, read two defensible ways, moves tracking error by 2.10
percentage points, eight times what the choice among four covariance estimators
moves the same quantity. The two are not close: a circular block bootstrap on
daily active returns puts the exemption effect at [1.84, 2.35] percentage
points and the estimator spread at [0.15, 0.40], and the first exceeds the
second in all 2{,}000 resamples.

We also report what the constraints buy, since the comparison invites the
question of whether the optimiser earns its complexity. On realised volatility
it does not: over this sample the optimised replicas are more volatile than the
heuristic one, not less. What the constraints deliver is control over where the
risk sits. The sector band binds at every review we examine, holding the
largest active industry weight at five per cent, while the unconstrained rule
drifts to between 32 and 48 per cent, growing monotonically across its reviews
as a single industry became both larger in the parent and more volatile. The
two sets of review dates do not coincide, so this is a contrast between a
binding constraint and its absence rather than a matched comparison; it is also
not a close call. Neither index produces a six-factor alpha significant
at the five per cent level, and betting-against-beta loadings between 8.7 and
14.0 absorb what intercept there is, which is consistent with the
leverage-constraint account of \citet{frazzini2014betting} and does not
distinguish it from the benchmarking account of \citet{baker2011benchmarks}.

Three boundaries on all of this. The parent universe is a proxy: the optimiser
works over the S\&P 500 where the real index uses a broader one, and that
substitution is the largest single contributor to the optimised tracking error
and cannot be repaired by anything else in the pipeline. The ordering between
specification and estimation is measured on one index pair over one sample
period, so it is a demonstration that the ordering can be this large, not a
claim that it is this large in general. And the minimum-holding rule makes the
programme mixed-integer; we handle it with the drop-and-resolve heuristic that
index providers themselves describe, and we do not measure the gap to an exact
solve. Section~\ref{sec:limits} sets these out in full.

Section~\ref{sec:rules} states both index rules from their primary documents and
identifies precisely what the optimised one leaves open.
Section~\ref{sec:model} builds the risk model and Section~\ref{sec:data} the
data behind it, including the three ways free price and fundamental data will
mislead a replication that does not defend against them.
Section~\ref{sec:design} gives the replication design.
Sections~\ref{sec:repro} to \ref{sec:anomaly} report what is reproducible,
whether the estimator matters, what the constraints buy, and whether any of it
is an anomaly. Section~\ref{sec:limits} and Section~\ref{sec:conclusion}
interpret and conclude.
